Niech $\C$ będzie kategorią.

\begin{definicja}
\emph{Diagramem produktowym} dla $A, B \in \C$ nazywamy obiekt $P$ wraz ze strzałkami
\[
\begin{tikzcd}
A & P \arrow[l, "p_1"'] \arrow[r, "p_2"] & B
\end{tikzcd}
\]
takimi, że dla każdego obiektu $X \in \C$ wraz ze strzałkami
\[
\begin{tikzcd}
A & X \arrow[l, "x_1"'] \arrow[r, "x_2"] & B
\end{tikzcd}
\]
istnieje jednoznaczne $u \colon X \to P$, które czyni diagram
\[
\begin{tikzcd}
& X \arrow[dl, "x_1"'] \arrow[dr, "x_2"] \arrow[d, "u" description, dashed] & \\
A & P \arrow[l, "p_1"] \arrow[r, "p_2"'] & B
\end{tikzcd}
\]
przemiennym. Piszemy $P = A \times B$ oraz $u = \langle x_1, x_2 \rangle$.
\end{definicja}

\begin{przyklad}
\leavevmode
\begin{enumerate}
\item W kategorii $\Set$ dla dwóch zbiorów $A, B \in \Set$ ich diagramem produktowym jest
\[
\begin{tikzcd}
A & A \times B \arrow[l, "\pi_1"'] \arrow[r, "\pi_2"] & B
\end{tikzcd}
\qquad A \times B = \{(a, b) \mid a \in A,\ b \in B\}.
\]

\item W omawianych wcześniej kategoriach algebr ($\Mon$, $\Grp$, $\Ring$ itd.) dla każdych dwóch obiektów $A, B$ istnieje produkt $A \times B$ zdefiniowany w naturalny sposób (np. dla $\Grp$ -- produkt grup).

\item Niech $\mathcal{P}$ będzie kategorią indukowaną przez $(P, \leq)$. Czym jest produkt $p \times q$ dla $p, q \in P$?
\[
\begin{tikzcd}
& p \times q \arrow[dl] \arrow[dr] & \\
p & & q
\end{tikzcd}
\]
\textbf{Obs.}
\begin{enumerate}
\item[a)] $p \times q \leq p$ oraz $p \times q \leq q$;
\item[b)] $\forall x$ t.\,że $x \leq p$, $x \leq q$ mamy $x \leq p \times q$.
\end{enumerate}
Widać, że $p \times q = p \wedge q$.
\end{enumerate}
\end{przyklad}

\begin{uwaga}
W kategorii indukowanej przez antyłańcuch $p \not\leq q$, $q \not\leq p$ nie ma produktu $p \times q$.
\end{uwaga}

\subsection*{Własności kategorii z produktem}

Niech $\C$ będzie kategorią, w której dla każdych dwóch obiektów istnieje ich diagram produktowy. Dodatkowo zakładamy, że mamy następujący diagram
\[
\begin{tikzcd}
A & A \times A' \arrow[l, "p_1"'] \arrow[r, "p_2"] & A' \\
B \arrow[u, "f"] & B \times B' \arrow[l, "q_1"] \arrow[r, "q_2"'] \arrow[u, dashed, "\exists! u"] & B' \arrow[u, "f'"']
\end{tikzcd}
\]
Oznaczamy $u = f \times f'$.

W podobny sposób możemy zdefiniować dla trzech obiektów $A, B, C$ ich produkt
\[
\begin{tikzcd}
& A \times B \times C \arrow[dl] \arrow[d] \arrow[dr] & \\
A & B & C
\end{tikzcd}
\quad (\ast)
\]

\begin{stwierdzenie}
Jeśli $\C$ jest kategorią z produktami binarnymi, to w $\C$ istnieją produkty ternarne $(\ast)$. Dla $A, B, C \in \C$ $(A \times B) \times C$ jest produktem ternarnym $A, B, C$.
\end{stwierdzenie}

\begin{definicja}
Mówimy, że kategoria $\C$ ma \emph{wszystkie skończone produkty}, jeśli
\begin{enumerate}
\item ma obiekt końcowy;
\item dla każdych $A, B \in \C$ istnieje nad nimi diagram produktowy.
\end{enumerate}
\end{definicja}

\subsection*{Więcej o dualności}

Niech $\C$ to:
\begin{itemize}
\item $\C_0$: $A, B, C$ -- obiekty;
\item $\C_1$: $f, g, h$ -- strzałki;
\item operacje: $\dom(f)$, $\cod(f)$, $\id_A$, $g \circ f$,
\end{itemize}
spełniające aksjomaty:
\[
\dom(\id_A) = A, \qquad \cod(\id_A) = A,
\]
\[
f \circ \id_{\dom(f)} = f, \qquad \id_{\cod(f)} \circ f = f,
\]
\[
\dom(g \circ f) = \dom(f), \qquad \cod(g \circ f) = \cod(g),
\]
\[
h \circ (g \circ f) = (h \circ g) \circ f.
\]

Niech $\Sigma$ będzie dowolnym zdaniem wyrażonym w podstawowym języku teorii kategorii. Możemy zdefiniować zdanie $\Sigma^{\op}$ dualne do $\Sigma$ zamieniając:
\begin{enumerate}
\item każde wystąpienie $f \circ g$ na $g \circ f$;
\item $\cod$ na $\dom$;
\item $\dom$ na $\cod$.
\end{enumerate}

Widać, że $\Sigma^{\op}$ też jest poprawnym zdaniem w TK. Zachodzi stąd: \emph{Dla każdego poprawnego zdania $\Sigma$, jeśli $\Sigma$ wynika z aksjomatów TK, to $\Sigma^{\op}$ też.}

\subsection*{Koprodukty}

Niech $\C$ będzie kategorią.

\begin{definicja}
Dla $A, B \in \C$ \emph{diagramem koproduktowym} nad $A, B$ nazywamy obiekt $Q$ wraz ze strzałkami
\[
\begin{tikzcd}
A \arrow[r, "q_1"] & Q & B \arrow[l, "q_2"']
\end{tikzcd}
\]
takimi, że dla każdego obiektu $Z$ oraz $A \xrightarrow{z_1} Z \xleftarrow{z_2} B$ istnieje jednoznaczne $u$ czyniące następujący diagram przemiennym
\[
\begin{tikzcd}
A \arrow[r, "q_1"] \arrow[dr, "z_1"'] & Q \arrow[d, dashed, "\exists! u"] & B \arrow[l, "q_2"'] \arrow[dl, "z_2"] \\
& Z &
\end{tikzcd}
\]
Oznaczenie: $u = [z_1, z_2]$, $Q = A + B$.
\end{definicja}

\begin{przyklad}
\leavevmode
\begin{enumerate}
\item Czym są koprodukty (jeśli istnieją) w $\Set$? Weźmy
\[
A + B = A \times \{1\} \cup B \times \{2\} \qquad (\text{rozłączna suma kopii}),
\]
oraz
\[
i_1 \colon A \to A + B, \quad a \mapsto (a, 1), \qquad i_2 \colon B \to A + B, \quad b \mapsto (b, 2).
\]
Wówczas
\[
\begin{tikzcd}
A \arrow[r, "i_1"] & A + B & B \arrow[l, "i_2"']
\end{tikzcd}
\]
jest diagramem koproduktowym nad $A, B$.

\item W kategorii $\mathcal{P}$ indukowanej przez poset $(P, \leq)$:
\[
\begin{tikzcd}
p \arrow[r] & p + q & q \arrow[l]
\end{tikzcd}
\]
Widać, że $p + q = p \vee q$.
\end{enumerate}
\end{przyklad}

Niech dany jest następujący diagram
\[
\begin{tikzcd}
A \arrow[r, "a_1"] \arrow[d, "f"'] & A + A' \arrow[d, dashed, "\exists! u"] & A' \arrow[l, "a_2"'] \arrow[d, "f'"] \\
B \arrow[r, "b_1"'] & B + B' & B' \arrow[l, "b_2"]
\end{tikzcd}
\]
Dla każdej pary $f \colon A \to B$, $f' \colon A' \to B'$ istnieje jednoznaczna strzałka
\[
f + f' \colon A + A' \to B + B'
\]
(o ile $\C$ jest kategorią z koproduktami binarnymi).

\begin{stwierdzenie}
Koprodukty są określone jednoznacznie z dokładnością do izomorfizmu.
\end{stwierdzenie}
\begin{proof}
Wynika z zasady dualności.
\end{proof}

\begin{definicja}
Mówimy, że $\C$ ma \emph{wszystkie skończone koprodukty}, jeśli:
\begin{enumerate}
\item $\C$ ma obiekt początkowy;
\item $\C$ ma wszystkie binarne koprodukty.
\end{enumerate}
\end{definicja}
